% Propuesta de inserción después de 07b_geometria_elipse.tex.
% RESULTADO_RECUPERADO: c52_06_01_incertidumbre.tex, líneas 10--91.
% Dependencias: c52_04_01_ley_areal.tex, líneas 137--198;
% c52_06_02_intercambio_constitutivo.tex, líneas 5--29.
% Explicitación local: dominio operatorio, prueba de covarianza,
% transporte recíproco y distinción dimensional posición--momento.
% Este fragmento no introduce una constante ni un generador adicional.

\subsection{Covarianza canónica y reciprocidad de la elipse de acción}
\label{sec:ii-area-covarianza}

La construcción precedente determina dos invariantes complementarios:
el producto de los semiejes fija el área de acción y su cociente conserva
la anisotropía del par angular. Precisamos ahora la realización canónica
en la que esos semiejes son dos desviaciones estándar. La estructura de
partida es la elipse ya obtenida de las representaciones angulares y de
acción de APP--TRIT--TPK; la realización operatoria se establece después
de esa construcción.

En este apartado escribimos \(\eta=\eta_{\rm el}
=\operatorname{artanh}(C^*/A)\), con \(|C^*/A|<1\) y
\(\hbar>0\). El cociente angular utiliza una misma unidad para sus dos
componentes. En las coordenadas equilibradas \((Q,P)\), ambas de
dimensión raíz de acción, definimos
\begin{equation}
 V_\eta:=\frac14
 \begin{pmatrix}a_S^2&0\\0&b_S^2\end{pmatrix}
 =\frac{\hbar}{2}
 \begin{pmatrix}e^\eta&0\\0&e^{-\eta}\end{pmatrix}.
 \label{eq:ii-covarianza-geometrica}
\end{equation}
Esta matriz positiva tiene entradas de dimensión acción. Su interpretación
como covarianza cuántica queda especificada por la realización siguiente.

\begin{proposition}[Realización gaussiana de la covarianza de acción]
\label{prop:ii-covarianza-gaussiana}
Considérese la representación canónica en \(L^2(\mathbb R,dQ)\), con
\(\widehat Qf(Q)=Qf(Q)\) y
\(\widehat Pf(Q)=-i\hbar f'(Q)\), inicialmente sobre el espacio de
Schwartz \(\mathcal S(\mathbb R)\). El estado normalizado
\begin{equation}
 \psi_\eta(Q)=\bigl(\pi\hbar e^\eta\bigr)^{-1/4}
 \exp\!\left(-\frac{Q^2}{2\hbar e^\eta}\right)
 \label{eq:ii-gaussiana-accion}
\end{equation}
tiene medias nulas y matriz de covarianza \(V_\eta\). En particular,
\begin{equation}
 (\Delta Q)^2=\frac{\hbar e^\eta}{2},\qquad
 (\Delta P)^2=\frac{\hbar e^{-\eta}}{2},\qquad
 \operatorname{Cov}(Q,P)=0,\qquad
 \det V_\eta=\frac{\hbar^2}{4}.
 \label{eq:ii-varianzas-accion}
\end{equation}
Esta realización satura la desigualdad de Robertson--Schrödinger.
\end{proposition}

\begin{proof}
Los operadores anteriores son esencialmente autoadjuntos sobre
\(\mathcal S(\mathbb R)\), que constituye un núcleo común invariante;
sobre este espacio se cumple
\([\widehat Q,\widehat P]=i\hbar I\). Para
\(\widehat Y=(\widehat Q,\widehat P)\), la covarianza simétrica es
\[
 V_{jk}=\frac12\left\langle
 (\widehat Y_j-\langle\widehat Y_j\rangle)
 (\widehat Y_k-\langle\widehat Y_k\rangle)
 +(\widehat Y_k-\langle\widehat Y_k\rangle)
 (\widehat Y_j-\langle\widehat Y_j\rangle)
 \right\rangle.
\]
Con \(s=\hbar e^\eta>0\), sea
\(I_s=\int_{\mathbb R}e^{-Q^2/s}\,dQ\). El producto de dos
integrales iguales, evaluado en coordenadas polares, da
\(I_s^2=2\pi\int_0^\infty r e^{-r^2/s}\,dr=\pi s\).
Además, la integral de
\(\frac{d}{dQ}(Qe^{-Q^2/s})\) es cero, por el decaimiento en ambos
extremos; de ahí
\(\int_{\mathbb R}Q^2e^{-Q^2/s}\,dQ=sI_s/2\).
Estas identidades y la paridad prueban la normalización, la media nula y
\(\langle\widehat Q^2\rangle=s/2\). Además,
\(\widehat P\psi_\eta=i e^{-\eta}Q\psi_\eta\).
Por tanto, \(\langle\widehat P\rangle=0\),
\(\|\widehat P\psi_\eta\|^2=\hbar e^{-\eta}/2\) y
\(\operatorname{Re}\langle\widehat Q\psi_\eta,
\widehat P\psi_\eta\rangle=0\).

Para cualquier estado normalizado de ese dominio, Cauchy--Schwarz aplicada
a \((\widehat Q-\langle\widehat Q\rangle)\psi\) y
\((\widehat P-\langle\widehat P\rangle)\psi\) da
\[
 (\Delta Q)^2(\Delta P)^2
 \geq \operatorname{Cov}(Q,P)^2
       +\frac14\left|\langle[\widehat Q,\widehat P]\rangle\right|^2
 =\operatorname{Cov}(Q,P)^2+\frac{\hbar^2}{4}.
\]
Las varianzas calculadas alcanzan la igualdad. Equivalentemente, el
operador
\[
 a_\eta=\frac{e^{-\eta/2}\widehat Q
                  +i e^{\eta/2}\widehat P}{\sqrt{2\hbar}}
\]
satisface \([a_\eta,a_\eta^\dagger]=I\) y
\(a_\eta\psi_\eta=0\), por sustitución directa.
\end{proof}

\begin{proposition}[Área, dispersión y transporte recíproco]
\label{prop:ii-area-covarianza-reciprocidad}
La elipse de acción es exactamente el subnivel de covarianza
\begin{equation}
 \mathcal E_\eta
 =\left\{y=(Q,P)^{\mathsf T}:y^{\mathsf T}V_\eta^{-1}y\leq4\right\}
 =\left\{\frac{Q^2}{a_S^2}+\frac{P^2}{b_S^2}\leq1\right\}.
 \label{eq:ii-elipse-covarianza}
\end{equation}
Su área es \(h=2\pi\hbar\), sus semiejes son
\(a_S=2\Delta Q\), \(b_S=2\Delta P\) y
\begin{equation}
 \frac{\Delta P}{\Delta Q}=e^{-\eta}
 =\frac{b_S}{a_S}=R_{\rm el}^2.
 \label{eq:ii-covarianza-forma}
\end{equation}
La rotación canónica \(J_2\) de la sección~\ref{iv:ellipse:section} transporta
\(V_\eta\) a \(V_{-\eta}\) y conserva la forma simpléctica.
\end{proposition}

\begin{proof}
La inversión de la matriz diagonal
\eqref{eq:ii-covarianza-geometrica} da
\eqref{eq:ii-elipse-covarianza}. Sus semiejes son las raíces positivas de
cuatro veces sus varianzas, y
\[
 \operatorname{Area}(\mathcal E_\eta)
 =4\pi\sqrt{\det V_\eta}=2\pi\hbar=h,
\]
en concordancia con la ley areal~\eqref{iv:ellipse:area-fase}.
La razón de las raíces positivas prueba
\eqref{eq:ii-covarianza-forma}.
Si \(M_\eta=\operatorname{diag}(e^{\eta/2},e^{-\eta/2})\), entonces
\[
 V_\eta=M_\eta V_0M_\eta^{\mathsf T},\qquad
 M_\eta^{\mathsf T}J_2M_\eta=J_2,\qquad
 J_2V_\eta J_2^{\mathsf T}=V_{-\eta}.
\]
Estas identidades se obtienen por multiplicación. En términos operatorios,
\((\widehat Q',\widehat P')=(-\widehat P,\widehat Q)\) conserva
\([\widehat Q',\widehat P']=i\hbar I\). La permutación
\(S_2:(Q,P)\mapsto(P,Q)\) también intercambia las varianzas, pero invierte
el signo del conmutador y de la forma simpléctica. Así, el transporte de
covarianza respeta la distinción entre reciprocidad canónica e inversión
de orientación ya establecida en la sección~\ref{iv:ellipse:section}.
\end{proof}

\paragraph{Posición, momento y transducción dimensional.}
Sea \(\lambda_{xp}>0\) un factor de la coordenatización dimensional, de dimensión
momento por unidad de longitud. La transformación
\[
 Q=\sqrt{\lambda_{xp}}\,x,\qquad
 P=\frac{p_x}{\sqrt{\lambda_{xp}}}
\]
realiza el par equilibrado en una posición \(x\) y su momento conjugado
\(p_x\); conserva \(dQ\wedge dP=dx\wedge dp_x\) y
\([\widehat x,\widehat p_x]=i\hbar I\). Por transformación de la
covarianza,
\begin{equation}
 V_\eta^{(x,p_x)}=\frac{\hbar}{2}
 \begin{pmatrix}e^\eta/\lambda_{xp}&0\\
                 0&\lambda_{xp}e^{-\eta}\end{pmatrix},\qquad
 \Delta x\,\Delta p_x=\frac{\hbar}{2},\qquad
 \frac{\Delta p_x}{\Delta x}=\lambda_{xp}e^{-\eta}.
 \label{eq:ii-covarianza-dimensiones}
\end{equation}
El cociente de dispersión adquiere aquí la dimensión momento por longitud.
En cambio, \eqref{eq:ii-covarianza-forma} pertenece a la coordenatización equilibrada.
Para la realización inductivo--capacitiva ya definida,
\eqref{hmtII:eq:ii-elipse-impedancia} aporta la identidad adimensional
\(Z_\eta/Z_*=e^{-\eta}=\Delta P/\Delta Q\).
Cada lectura utiliza su transducción declarada y conserva el mismo
determinante de acción.

\paragraph{Energía del contorno y energía media del estado.}
La realización inductivo--capacitiva tiene el Hamiltoniano
\[
 H_\eta(Q,P)=\frac{\omega}{2}
 \left(e^{-\eta}Q^2+e^\eta P^2\right),\qquad \omega>0.
\]
Sobre \(\partial\mathcal E_\eta\), su valor es
\(H_\eta=\hbar\omega\), por las ecuaciones de los semiejes. En la
realización cuántica de la proposición~\ref{prop:ii-covarianza-gaussiana},
\begin{equation}
 \langle\widehat H_\eta\rangle_{\psi_\eta}
 =\frac{\omega}{2}
 \left(e^{-\eta}\frac{\hbar e^\eta}{2}
       +e^\eta\frac{\hbar e^{-\eta}}{2}\right)
 =\frac{\hbar\omega}{2}.
 \label{eq:ii-energia-contorno-covarianza}
\end{equation}
El contorno de dos desviaciones estándar tiene, por tanto, el doble de la
energía media de este estado gaussiano. La ley areal fija la normalización
de acción; el par angular fija su distribución entre las dos coordenadas;
la representación canónica y el estado explícito realizan esa distribución
como covarianza. La saturación de incertidumbre queda así acreditada en
la realización especificada por sus operadores y su estado.
